% ============================================================
%  MM-GTUNets 课堂汇报 · 主题
%  蓝白主题，主色取自西工大校徽钴蓝（#004088）
%  颜色即语义：影像 = 钴蓝，非影像 = 琥珀，融合 / 本文方法 = 藏青
%  钴蓝与琥珀只表示模态；其余强调一律用藏青系；红色只用于“惩罚”
% ============================================================
\usefonttheme{professionalfonts}
\usepackage{fontspec}
\usepackage{xeCJK}
\usepackage{amsmath}
\usepackage{unicode-math}
\usepackage{tikz}
\usetikzlibrary{arrows.meta,calc,positioning,fit,backgrounds,decorations.pathreplacing}

% ---------- 字体（均为 SIL OFL 开源字体，文件位于 fonts/） ----------
%  中文：思源黑体 Source Han Sans CN
%  西文：Fira Sans（Metropolis 主题的标准字体）
%  数学：Fira Math（与 Fira Sans 配套设计的无衬线数学字体）
\defaultfontfeatures{Path=fonts/}
\setCJKmainfont{SourceHanSansCN-Regular.otf}[BoldFont=SourceHanSansCN-Bold.otf]
\setCJKsansfont{SourceHanSansCN-Regular.otf}[BoldFont=SourceHanSansCN-Bold.otf]
\setmainfont{FiraSans-Regular.otf}[BoldFont=FiraSans-SemiBold.otf,ItalicFont=FiraSans-Italic.otf,BoldItalicFont=FiraSans-SemiBoldItalic.otf]
\setsansfont{FiraSans-Regular.otf}[BoldFont=FiraSans-SemiBold.otf,ItalicFont=FiraSans-Italic.otf,BoldItalicFont=FiraSans-SemiBoldItalic.otf]
\setmathfont{FiraMath-Regular.otf}[Path=fonts/]
\xeCJKsetup{CJKecglue={\hskip 0.16em plus 0.08em}}

% ---------- 颜色 ----------
\definecolor{Ink}{HTML}{1E2532}
\definecolor{Navy}{HTML}{16294A}
\definecolor{NavyPale}{HTML}{E8EEF7}
\definecolor{NavyMid}{HTML}{4A5D85}
\definecolor{Muted}{HTML}{667183}
\definecolor{Faint}{HTML}{A3ACBA}
\definecolor{Line}{HTML}{DDE3EA}
\definecolor{Pale}{HTML}{F4F7FB}
\definecolor{Img}{HTML}{2B64AE}
\definecolor{ImgPale}{HTML}{E6EEF8}
\definecolor{Non}{HTML}{C47A2C}
\definecolor{NonPale}{HTML}{FAEEE1}
\definecolor{Case}{HTML}{34405E}
\definecolor{Ctrl}{HTML}{C4CEDD}
\definecolor{Rew}{HTML}{16294A}
\definecolor{Pen}{HTML}{C0413B}

% ---------- 字号 ----------
\newcommand{\fbody}{\fontsize{10}{15}\selectfont}
\newcommand{\fsmall}{\fontsize{8.6}{12.4}\selectfont}
\newcommand{\fnote}{\fontsize{7.8}{11}\selectfont}
\newcommand{\ftiny}{\fontsize{7.6}{9.4}\selectfont}
\newcommand{\fhead}{\fontsize{10.5}{14}\selectfont\bfseries}
\newcommand{\fbig}{\fontsize{20}{22}\selectfont\bfseries}

% ---------- Beamer 外观 ----------
\setbeamercolor{normal text}{fg=Ink,bg=white}
\setbeamercolor{background canvas}{bg=white}
\setbeamercolor{structure}{fg=Navy}
\setbeamertemplate{navigation symbols}{}
\setbeamersize{text margin left=0.9cm,text margin right=0.9cm}

\newcommand{\kickertext}{}
\newcommand{\kick}[1]{\renewcommand{\kickertext}{#1}}
\newcounter{partno}
\newcommand{\setpart}[1]{\setcounter{partno}{#1}}
\newcommand{\totalpages}{15}

% 方法部分的小导航：构图 › 编码 › 融合，当前环节实心高亮（0 = 不显示）
\newcounter{navstage}
\newcommand{\setstage}[1]{\setcounter{navstage}{#1}}
\tikzset{
  pillon/.style={rounded corners=5pt,fill=Navy,draw=Navy,line width=.6pt,text=white,inner xsep=6pt,inner ysep=2.6pt,
                 font=\fontsize{7.8}{9}\selectfont\bfseries},
  pilloff/.style={rounded corners=5pt,fill=white,draw=Line,line width=.6pt,text=Muted,inner xsep=6pt,inner ysep=2.6pt,
                  font=\fontsize{7.8}{9}\selectfont},
}
\newcommand{\navpill}[3]{%  #1 环节编号  #2 文字  #3 右端坐标
  \ifnum\value{navstage}=#1\relax\def\pillstyle{pillon}\else\def\pillstyle{pilloff}\fi
  \node[\pillstyle,anchor=east] (pill#1) at #3 {#2};}

\setbeamertemplate{frametitle}{%
  \begin{tikzpicture}[remember picture,overlay]
    \node[anchor=north west,inner sep=0,font=\fontsize{7.6}{9}\selectfont\bfseries,text=NavyMid]
      at ([xshift=.9cm,yshift=-.42cm]current page.north west) {\kickertext};
    \node[anchor=north west,inner sep=0,font=\fontsize{17}{20}\selectfont\bfseries,text=Navy]
      at ([xshift=.9cm,yshift=-.74cm]current page.north west) {\insertframetitle};
    \node[anchor=north east,inner sep=0]
      at ([xshift=-.9cm,yshift=-.42cm]current page.north east) {\includegraphics[height=.68cm]{assets/nwpu-crest.png}};
    \ifnum\value{navstage}>0\relax
      \navpill{3}{融合}{([xshift=-2.05cm,yshift=-.78cm]current page.north east)}
      \navpill{2}{编码}{([xshift=-.3cm]pill3.west)}
      \navpill{1}{构图}{([xshift=-.3cm]pill2.west)}
      \foreach \a/\b in {1/2,2/3}
        \draw[Faint,line width=.6pt] ([xshift=.05cm]pill\a.east) -- ([xshift=-.05cm]pill\b.west);
    \fi
  \end{tikzpicture}}

\newcommand{\partlabel}[2]{%
  \ifnum\value{partno}=#1\relax
    {\color{Navy}\bfseries #2}%
  \else
    {\color{Faint}#2}%
  \fi}
\setbeamertemplate{footline}{%
  \begin{tikzpicture}[remember picture,overlay]
    \node[anchor=south west,inner sep=0,font=\ftiny]
      at ([xshift=.9cm,yshift=.22cm]current page.south west)
      {\partlabel{1}{问题}\hspace{.5em}\textcolor{Line}{/}\hspace{.5em}%
       \partlabel{2}{方法}\hspace{.5em}\textcolor{Line}{/}\hspace{.5em}%
       \partlabel{3}{实验}\hspace{.5em}\textcolor{Line}{/}\hspace{.5em}%
       \partlabel{4}{总结}};
    \node[anchor=south east,inner sep=0,font=\ftiny,text=Faint]
      at ([xshift=-.9cm,yshift=.22cm]current page.south east)
      {\textcolor{Muted}{\insertframenumber}\,/\,\totalpages};
  \end{tikzpicture}}

% ---------- 画布：内容区左下角为原点，宽 14.2，高 6.55（单位 cm） ----------
\newenvironment{canvas}{%
  \begin{tikzpicture}[remember picture,overlay,x=1cm,y=1cm,
     shift={($(current page.south west)+(.9cm,.6cm)$)}]}{%
  \end{tikzpicture}}

% ---------- 绘图样式 ----------
\tikzset{
  >={Stealth[length=2.1mm,width=1.6mm]},
  txt/.style={anchor=north west,inner sep=0,align=left,font=\fbody,text=Ink},
  head/.style={anchor=north west,inner sep=0,align=left,font=\fhead,text=Navy},
  note/.style={anchor=north west,inner sep=0,align=left,font=\fnote,text=Muted},
  box/.style={rounded corners=2.2pt,draw=Line,line width=.6pt,fill=white,inner xsep=5pt,inner ysep=4pt,
              align=center,font=\fsmall,text=Ink},
  imgbox/.style={box,draw=Img!70,fill=ImgPale},
  nonbox/.style={box,draw=Non!70,fill=NonPale},
  navbox/.style={box,draw=Navy!60,fill=NavyPale},
  arr/.style={->,line width=.8pt,draw=Faint},
  imgarr/.style={->,line width=1pt,draw=Img},
  nonarr/.style={->,line width=1pt,draw=Non},
  navarr/.style={->,line width=1pt,draw=Navy},
  gnode/.style={circle,minimum size=.36cm,inner sep=0,line width=.7pt},
  case/.style={gnode,fill=Case,draw=Case},
  ctrl/.style={gnode,fill=Ctrl,draw=Ctrl!80!Case},
  unk/.style={gnode,fill=white,draw=Faint,dashed},
  gedge/.style={draw=Navy,line width=#1},
  rule/.style={draw=Line,line width=.6pt},
}

% 页面底部的结论条
\newcommand{\takeaway}[1]{%
  \draw[rule] (0,.86) -- (14.2,.86);
  \fill[Navy] (0,.2) rectangle (.07,.54);
  \node[anchor=west,inner sep=0,font=\fbody\bfseries,text=Navy] at (.28,.37) {#1};}

\newcommand{\stepno}[1]{\tikz[baseline=-.6ex]\node[circle,fill=Navy,text=white,inner sep=0,minimum size=.3cm,font=\fontsize{6.5}{7}\selectfont\bfseries]{#1};}
\newcommand{\img}[1]{\textcolor{Img}{#1}}
\newcommand{\non}[1]{\textcolor{Non}{#1}}
\newcommand{\nav}[1]{\textcolor{Navy}{\bfseries #1}}
\newcommand{\muted}[1]{\textcolor{Muted}{#1}}
